\documentclass[12 pt, letterpaper]{hmcpset}
\usepackage{hw-preamble}

\name{Jayden Thadani}
\class{MATH 168 --- Algorithms}
\assignment{Homework \#10}
\duedate{2nd December 2025}

\begin{document}

\begin{problem}[Problem 1: Subset some?]
	Consider the optimization version of the Subset Sum problem: 
\begin{itemize}
    \item \textbf{Input:} A set $S$ of $n$ positive integers $\{v_1, \dots v_n\}$ and a target value $W$.
    \item \textbf{Output:} A subset $T \subset S$ with the largest possible sum not exceeding $W$.
\end{itemize}
For instance, if I had the set $S = \{2, 5, 11, 9\}$ and my target was $W = 15$, then the subset I
would choose is $T = \{5, 9\}$ with a sum of 14. (For this problem, feel free to
assume that $S$ won't contain any elements larger than $W$, as those would be easy
to identify and remove in linear time.)

\begin{enumerate}
    \item First, consider the following greedy algorithm proposed by Dr. D.
    Ception: 
    \begin{enumerate}[label = (\arabic*)]
        \item Let $T$ be an empty set.
        \item For $i = 1, 2, \dots, n$:
        \begin{itemize}
            \item Add $v_i$ to $T$ if this won't cause $T$ to exceed the target value $W$.
        \end{itemize}
        \item Return $T$.
    \end{enumerate}
    Dr. D. Ception claims this algorithm is (1/2)-optimal, i.e. that the value of the sum of $T$ found by the algorithm is at least half the optimal value
    one could achieve. 
    
    In a sentence or two, give a counterexample to the claim that the algorithm is (1/2)-optimal: that is, specify an instance on which the algorithm might fail to achieve a (1/2)-approximation of the optimal solution and explain how the algorithm could fail.
    
    \item In fact, the algorithm doesn't achieve any constant approximation factor! Explain how, for any integer constant $c$, you can modify your counterexample to get an instance on which the algorithm might fail to achieve a $(1/c)$-approximation. Briefly explain why the algorithm fails. (Please stick to positive integers.)

    \item Dr. D. Ception's algorithm can be modified to produce a (1/2)-optimal solution---it works \emph{if} the elements $v_1, v_2, \dots, v_n$ are considered in a certain order! Please
    \begin{itemize}
        \item identify the ordering of elements that will make D. Ception's algorithm return a (1/2)-approximation,
        \item \emph{prove} that an algorithm that begins by sorting the elements in this particular order and then using D. Ception's procedure is (1/2)-optimal, and
        \item state the runtime of the new algorithm.
    \end{itemize}
\end{enumerate}
\end{problem}

\begin{solution}
	\begin{enumerate}
		\item Let $S = \{ 1, W \}$ (in that order) for target value $W > 2$. Then D. Ception's algorithm gives subset $T = \{ 1 \}$ which has sum $1$. The actual optimal sum is $W$ by choosing $T = \{ W \}$. Thus, D. Ception's algorithm achieves only a $1 / W$ approximation with $1 / W < 1 / 2$. That is, it is not $1 / 2$-optimal.

		\item By choosing $W > c$ in the example above we get an algorithm that can only achieve a $1 / W$ approximation. Since $1 / W < 1 / c$, it is not even $1 / c$-optimal.

		\item Sort $v_{1}, \dots, v_{n}$ in descending order. That is, sort them so that $v_{i} \ge v_{i + 1}$ for each $i$. Then run D. Ception's algorithm on this sorted list.

			We will show that running D. Ception's algorithm on the sorted $S$ cannot give a set $T$ with $\Sigma(T) < W / 2$ unless $T$ itself is optimal. Since any optimal $U$ must have $\Sigma(U) \le W$, this gives $\Sigma(T) \ge \Sigma(U) / 2$, and thus, the greedy algorithm is $1/2$-optimal.

			Suppose $S$ contains an element $s \ge W / 2$. Then since $S$ contains no elements greater than $W$, the first element in $T$ is at least $W / 2$. Thus, $\Sigma(T) \ge W / 2 \ge \Sigma(U) / 2$.

			Now suppose $S$ contains only elements less than $W / 2$. Suppose some $T' \subseteq S$ has $\Sigma(T') < W / 2$. If $T'$ is not optimal, then there is some $s \notin T'$. Since $\Sigma(T') < W / 2$ and $s < W / 2$ we have $\Sigma(T \cup \{ s \}) \le W$. Thus, $T'$ cannot come from D. Ception's greedy algorithm --- adding any element of $S$ to $T$ obtained from the greedy algorithm would increase the sum beyond $W$. Thus, if $T$ comes from the greedy algorithm, either $\Sigma(T)$ is optimal, or $\Sigma(T) \ge W / 2$. In either case, for optimal $U$, $\Sigma(T) \ge \Sigma(U) / 2$.

Sorting $S$ by mergesort runs in $O(n \log n)$. Then running D. Ception's algorithm would involve $n$-many $O(1)$ checks (with possible inserts), and thus would run in $O(n)$. In total, this makes for $O(n \log n) + O(n) = O(n \log n)$ runtime.
	\end{enumerate}
\end{solution}

\pagebreak

\begin{problem}[Problem 2: Madcap map mark-up]
As a budding fantasy author, you've recently submitted the final manuscript of your soon-to-be-printed debut novel to your publisher. All that remains is to work out the practical details for the hardback version. You've drafted a beautiful black-and-white sketch of the world map that will be printed inside the front cover. (Every good fantasy novel should have a map at the front.) 

The printers will handle the map coloring for you using a default palette containing 4 colors. Unfortunately, they've never heard of the four-color theorem! Instead, they color every region on the map by choosing a color at random. This is unfortunate for your map, but how bad are things likely to look?

\vspace{0.5cm}

Your map consists of a set $R$ of regions in the plane. Some regions border each other. You may assume all borders have positive length (that is, don't worry about the situation in which many regions meet at a single point).

\begin{enumerate}
    \item We say that a border between two regions $r_i$ and $r_j$ is \emph{satisfied} if the two regions are different colors. Over the entire graph $G$, what is the expected fraction of satisfied borders when regions are colored uniformly at random?
    
    \item  Prove that, for \emph{any} map, randomly coloring the regions with 4 colors will satisfy \emph{at least} $1/2$ of the edges with probability \emph{at least} $1/2$.

    \emph{Hint:} This question requires a standard kind of probabilistic argument using the expected value calculation from the previous part. If that's something you haven't seen before and you find this difficult to prove, wait until we've had our lecture on randomized algorithms.
    
    \item The citizens of your fantasy world are enthusiasts of geometrical structure (aren't we all?) and as such have chosen to tile the landscape with hexagonal provinces. This results in a hexagonal lattice in which every province has 6 neighbors. (See figure for an example hexagonal lattice.)
    
    We say that a province is \emph{satisfied} if \emph{all} adjacent borders are satisfied. What is the chance that a random assignment of 4 colors satisfies a province with 6 neighbors? 
    
        \item After several frustrating conversations with the printers, you're considering paying them to use additional colors (they'll still assign these additional colors at random).

    What is the minimum number of colors required such that a province with 6 neighbors is satisfied with probability at least 50\%?

    \item That's more colors than you can afford! In despair, you turn to your longtime frenemy Dr. D. Ception, who is always quick to come up with a greedy algorithm. He recommends the following greedy coloring algorithm, which is simple enough that the printers should be able to follow it:
    \begin{enumerate}[label = (\arabic*)]
        \item Pick \emph{any} uncolored province.
        \item Color that vertex \emph{any} previously used color not already taken by one of its neighbors. If the vertex is adjacent to \emph{all} previously used colors, color it using a new color.
        \item Repeat until all vertices are colored.
    \end{enumerate}
    How many colors will this algorithm use to color a lattice of hexagonal province, in the worst case?
\end{enumerate}
\end{problem}

\begin{solution}
	\begin{enumerate}
		\item First we calculate the expectation of the indicator $1_{e}$ function that indicates whether a particular edge $e$ is satisfied or not. $e$ is satisfied if the two regions that it borders are coloured differently. Among the $16$ $4$-colouring possibilities for the two regions, $4$ of them have the two regions coloured the same. In these cases, the value of $1_{e}$ is $0$. In the remaining, the value of $1_{e}$ is $1$. Thus,
			\[
				\mathbb{E}(1_{e}) = 1 \times 12 / 16 + 0 \times 4 / 12 = 3/4.
			\]

			The expected fraction of satisfied borders is the expected value of the sum of all $1_{e}$ for all edges $e \in E(G)$. By linearity of expectation this is
			\[
				\mathbb{E}(X) = \sum_{e \in E(G)} \mathbb{E}(1_{e}) = (3 / 4) \# E(G).
			\]
			where $X = \#\text{satisfied edges}$. That is, we can expect $3 / 4$th of the edges to be satisfied.

		\item Let $X$ be the number of satisfied edges and let $m$ be the total number of edges. Then by the definition of expected value
			\begin{align*}
				3m/4 & = \sum_{i = 1}^{m} \mathbb{P}(X = i) i \\
				     & \le \sum_{i = 1}^{m / 2 - 1} \mathbb{P}(X = i) i + \sum_{i = m / 2}^{m} \mathbb{P}(X = i) i \\
				     & \le m/2 \sum_{i = 1}^{m / 2 - 1} \mathbb{P}(X = i) + m \sum_{i = m / 2}^{m} \mathbb{P}(X = i).
			\end{align*}
			Note that $\mathbb{P}(X \ge m / 2) = \sum_{i = m / 2}^{m} \mathbb{P}(X = i)$. Call this probability $p$. Then we can write
			\[
				3m / 4 \le (1 - p) m / 2 + pm = (1 + p)m/2.
			\]
			Then
			\[
				\frac{3m}{4} \frac{2}{m} \le 1 + p
			\]
			and finally
			\[
				1 / 2 \le p.
			\]

		\item This is the probability that the province itself is not coloured the same as any of its six neighbours. For this, each of the $6$ neighbours must be coloured one of the other $3$ colours, which happens independently for each with probability $3 / 4$. That is, the probability is $(3 / 4)^{6}$.

		\item By our previous answer, we want the minimal $k$ such that $(\frac{k - 1}{k})^{6} > 1/ 2$. Since $(8 / 9)^{6} < 1/2$ and $(9 / 10)^{6} > 1/2$, the minimal $k$ is $k = 10$.

		\item Since a hexagon in a hexagonal lattice has at most $6$ bordering hexagons, the maximum number of colours is bounded above by $7$. That is, if there are $7$ colours used, then any province borders at most $6$ provinces, and thus, there is always a previously used colour that the province is not adjacent to.

			Choosing a bad order of provinces to apply the greedy algorithm can actually force us to use $7$ colours. Essentially, we colour hexagons so that each hexagon surrounding a central ``capital'' hexagon must have a different colour. Then the capital must be coloured a distinct $7$th colour. This can be achieved by the colouring scheme below.

			\begin{figure}[H]
				\centering
				\includegraphics[width = 0.5 \textwidth]{figures/P2 e colouring.pdf}
				\caption{Colour the hexagons one of $1, \dots, 7$, in the order indicated by the pencil, starting from the circle, and ending at the hexagon coloured $7$}.
			\end{figure}
	\end{enumerate}
\end{solution}

\begin{problem}[Problem 2 : Madcap map markup]
	How many colors will this algorithm use to color a lattice of hexagonal provinces, in the worst case?
\begin{enumerate}
	\item[(f)] You've had enough of these shenanigans from the printers---you'll just color the map yourself. What is the smallest number $k$ of colors required to $k$-color an infinite lattice of hexagonal provinces? (Feel free to ``prove" your answer by drawing a picture.)

	\item[(g)] (10/40 points; a bit harder.) As it turns out, $4$-coloring is in $\mathsf{P}$. However, just for fun, you've decided to try out a variant of the exponential-time algorithm for 3-SAT we saw in our lecture on randomized algorithms. Given a planar map with $n$ regions, your algorithm does the following:

    \begin{enumerate}[label = (\arabic*)]
        \item Begin with a random coloring: assign each region one of the four ``colors'' $\{0, 1, 2, 3\}$ independently and uniformly at random.
        \item If all borders are satisfied, halt and return the coloring. If not, choose any border that is not satisfied; i.e., the regions on each side have the same coloring. Call these regions $r_i$ and $r_j$.
        \item Pick one of $r_i$ and $r_j$ at random, then assign it a \emph{new} color by choosing uniformly at random from the three available possibilities.
        \item Repeat steps (2) and (3) $3n/4$ times. If we don't find a valid coloring during that time, return FAIL.
    \end{enumerate}

    You also have a \emph{reference coloring} $C: \{0, 1, 2, ... n-1\} \rightarrow \{0, 1, 2, 3\}$ that assigns every region a color from the set $\{0, 1, 2, 3\}$ and satisfies all borders (thus creating a valid 4-coloring). After running step (1) of your algorithm, you observe that the initial random coloring matches $C$ on exactly $n/4$ regions. 

    Using the number of regions that match $C$ as a helpful reference metric, prove an (exponential) lower bound on the probability that the algorithm returns a valid 4-coloring. Write your result in the form $\alpha^{-n}$, where $\alpha$ should be a constant greater than $1$. How could you use this bound to create a (very slow) randomized algorithm for 4-coloring? Would this beat the naive $4^n$ brute-force algorithm?

    \emph{Hint:} No need to do fancy calculations here (although you can try if you want). It suffices to use the method we used for 3SAT to lower-bound the chance that we guess a coloring ``perfectly''.
\end{enumerate}
\end{problem}

\begin{solution}
	\begin{enumerate}
		\item[(f)] Since the hexagonal lattice contains $K_{3}$ --- there are triples of provinces, each adjacent to all others --- the lattice cannot be $2$-colourable. We show that it is in fact $3$-colourable, and thus, the minimal $k$ is $k = 3$. The $3$-colouring is obtained by $3$-colouring each set of $7$ hexagons in the manner below.
			\begin{figure}[H]
				\centering
				\includegraphics[width = 0.5 \textwidth]{figures/P2 f 3-colouring.pdf}
				\caption{A $3$-colouring of the hexagonal lattice}
			\end{figure}

		\item[(g)] Our random colouring agrees with a correct colouring $C$ in $t = n/4$ regions. If  $r_{i}, r_{j}$ is a pair of bordering regions coloured wrong, then at least one of them disagrees with $C$. By randomly picking one and then randomly changing its colour to one of the other $3$ colours, we have a $1 / 6$th chance of changing the correct one to the correct thing. That is, at each step $t$ increases to $t + 1$ with probability at least $1 / 6$. In order to get to $t = n$ from $t = n / 4$ we need to make $3n / 4$ correct guesses in a row with probability $(1 / 6)^{3n/4} = (6^{3/4})^{-n} \ge 3.833^{-n}$.

			Thus, trying our procedure $O(3.9^{n})$ times, we expect to get an exact $4$-colouring. Since our procedure is sub-exponential the total runtime of this algorithm is better than the bruteforce $O(4^{n})$ algorithm.
	\end{enumerate}
\end{solution}

\begin{problem}[Problem 3: The greediest algorithm]
	Welcome to the world-famous game show \emph{Cash Grab}! In this game, the contestant is presented with an $n \times n$ game board. Every square on the game board is associated with some non-negative amount of money, and no two squares have the same amount. The game proceeds as follows:
\begin{enumerate}
    \item The contestant selects a square $s$ on the game board and receives the amount of money associated with that square.
    \item All squares that share an edge with $s$ have their money \emph{removed} from the game. The contestant can no longer choose these squares or receive money from them. (If the square has already been emptied, nothing happens to it.) 
\end{enumerate}
Play continues until there's no more cash left on the board. 

\vspace{0.5cm}

One popular greedy strategy is to repeatedly choose the most valuable square remaining until there is no money left. In this question, you'll analyze: how good is this strategy?
\begin{enumerate}
    \item A set of game board squares is \emph{valid} if no two squares share an edge (this implies that a contestant can choose all of them). Let $OPT$ be any valid set of squares that maximizes total winnings, and let $GREEDY$ be the valid set of squares chosen by a contestant who follows the greedy strategy. 
    
    In a few sentences, show that for every square $s \in OPT$, either (1) $s \in GREEDY$ or (2) $s$ is adjacent to a square $t \in GREEDY$ such that $t$ has a \emph{higher} value than $s$.
    
    \item Show that a contestant following the greedy algorithm always receives at least 1/4 of the optimal amount that they could possibly get. That is, prove that the greedy algorithm is $(1/4)$-optimal.

    \item In fact, there is a very simple $O(n^2)$-time strategy that achieves a (1/2)-approximation. What is it? Justify your answer.
\end{enumerate}
\end{problem}

\begin{solution}
	\begin{enumerate}
		\item If $s \in GREEDY$, we are done. Suppose $s \notin GREEDY$. If $s$ has greater value than all of its neighbours, then none of its neighbours would be picked over it, leaving it with positive value until it is picked over its neighbours. Then $s$ would be picked in $GREEDY$. Thus, $s \notin GREEDY$ must have some neighbour $t$ with greater value than it.

		\item Let $OPT$ give the set $S$ and $GREEDY$ give the set $T$. Then $\Sigma(S) - \Sigma(T) = \Sigma(S \setminus T) - \Sigma(T \setminus S)$. That is, the difference is the difference between the values chosen in only one of them since the values of the squares chosen in both just cancel out. For every $s \in S \setminus T$, there is some $t \in T$ adjacent to $s$ and with greater value. In fact, since $s$ is adjacent to $t$, $t$ is not in $OPT$ and so  $t \in T \setminus S$. Since each square on the grid has at most four neighbours, each $t$ is adjacent to at most four $s \in S \setminus T$ (all with value less than it). Then $\Sigma(S \setminus T) \le 4 \Sigma(T \setminus S)$. Thus,
			\begin{align*}
				\Sigma(S) - \Sigma(T) & = \Sigma(S \setminus T) - \Sigma(T \setminus S) \\
						      & \le 4\Sigma(T \setminus S) - \Sigma(T \setminus S) \\
						      & =  3 \Sigma(T \setminus S) \\
						      & \le 3 \Sigma(T).
			\end{align*}
			Then it is easy to see that $\Sigma(S) - \Sigma(T) \le 3 \Sigma(T)$ implies $\Sigma(S) \le 4 \Sigma(T)$ and thus, $\Sigma(T) \ge \Sigma(S) / 4$ --- the greedy algorithm is at least $1 / 4$-optimal.

		\item We first describe this algorithm. Colour the grid black and white like a chessboard. Note that since no two squares of the same colour share an edge, the set of all black squares $B$ and the set of all white squares $W$ are both valid. We claim that choosing the larger of the two sums is $1 / 2$-optimal.

			We prove $1/2$-optimality as follows. Since every square is black or white, $\Sigma(B) + \Sigma(W)$ is the sum of the values of all squares. Since the optimal value is certainly bounded by the sum of values of all squares, we have $\Sigma(B) + \Sigma(W) \ge \Sigma(OPT)$. Suppose without loss of generality that $\Sigma(B) \ge \Sigma(W)$. Then clearly $2 \Sigma(B) \ge \Sigma(B) + \Sigma(W) \ge \Sigma(OPT)$. That is, $\Sigma(B) \ge \Sigma(OPT) / 2$. Thus, choosing $\max \{ \Sigma(B), \Sigma(W) \}$ is $1 / 2$-optimal.

			To compute the sum of all black squares involves $n^{2} / 2 = O(n^{2})$ summations and so does the sum of all white squares. Comparing the two sums can be done in constant time so the algorithm runs in $O(n^{2} + 1) = O(n^{2})$ time.
	\end{enumerate}
\end{solution}

\pagebreak

\begin{problem}[Problem 4: Multicommodity flow is $\mathsf{NP}$-hard!]
	We know how to solve the network flow problem in polynomial time.  A generalization of
		network flow is called \emph{multicommodity network flow} and it goes like this:
		We are given a flow network except now there are multiple source/sink pairs in the network.
		That is, there are some number $k$ of sources $s_1, \dots, s_k$ and the same number of sinks $t_1, \dots, t_k$.
		Each source $s_i$ pushes a different kind of ``fluid'' or ``commodity'' to its corresponding sink $t_i$.  However, the edges of the network are shared by these flows and each edge still has a single
		integer capacity.  For example, an edge with capacity 5 might accommodate 2 units of flow from $s_1$ to $t_1$ and 3 units of flow from $s_2$ to $t_2$.  Now we are interested in finding a set of flows, one for each source/destination pair.  Of course, we are still subject to capacity constraints and conservation of flow (i.e. for any given commodity, the amount of flow for that commodity entering a given vertex is equal to the amount of flow for that commodity leaving that vertex).  In addition, we require the flow of each commodity on an edge to be an integer.  Amazingly, this generalization of the network flow problem is NP-hard!
		 
	   In particular, the decision problem that we will consider here is this:  Given is a flow network with $k$ commodities and $k$ source/sink pairs, $(s_i, t_i)$ (one for each commodity), and a demand $d_i$ for each sink $t_i$.  The question is:  ``Does there exist a valid set of flows, one for each commodity, such that each sink receives $d_i$ of commodity $i$?''
		
		Prove this problem is NP-complete by reduction from 3SAT.
\end{problem}

\end{document}
